\chapter{Conclusions and Recommendations}
\label{chap:conclusion}

This chapter presents the conclusions supported by the analytical, CAD, and
bracket-study evidence. It also states the limitations of the work and the
validation required before any candidate can progress toward manufacture or
service.

\section{Conclusions}
\label{sec:overall-conclusion}

The study developed and evaluated an auditable, resource-conscious workflow
for the parameterized mass optimization of three marine structural component
archetypes. The same overall process represented a bracket, padeye, and
stabilizer while allowing each component to retain its own variables,
material, load, analytical formulation, CAD generator, and boundary roles.
This supports the use of a shared part-aware architecture within the reported
scope.

Seeded Latin-hypercube exploration and three bounded Nelder--Mead starts
identified lower-mass candidates that satisfied the implemented analytical
factor-of-safety and displacement criteria. The selected screening reductions
were 72.9\% for the bracket, 59.9\% for the padeye, and 27.2\% for the
stabilizer relative to their configured baselines. Multiple starts improved
the best feasible mass for every part. These values describe the best
feasible designs found within the declared parameterized searches; they do not
prove global optimality or service-ready weight savings.

Native Fusion generation remained usable across the seeded CAD campaigns. The
bracket and padeye completed all 21 baseline-plus-sample cases. The stabilizer
created all 21 artifact sets, although the outer process timed out while
reading the final response. Analytical and native-CAD mass differed by no more
than 3.8\% for the six baseline and finalist comparisons. This agreement
supports the volume models for early mass ranking, but it does not validate
stress or displacement.

The two bracket Static Stress studies provide the clearest lesson about
evidence quality. Their images identify the geometry, mounting-hole supports,
rib-face loading, and qualitative response regions. Study~2 also records the
mesh and boundary setup more clearly than Study~1. Nevertheless, Study~2 used
1.472~N instead of the intended 1472~N, while Study~1 did not provide a
complete traceable load definition and reported only a 1.298~N support
reaction. The numerical stress, factor of safety, and displacement values from
both studies were therefore rejected.

The principal outcome is a traceable method for reducing a large design set,
checking native geometry and mass, and preparing a small candidate set for
disciplined validation. The study does not establish structural
qualification, fabrication readiness, or service safety.

\section{Achievement of the Objectives}
\label{sec:aim-achievement}

The aim was achieved at the level of workflow development, analytical
screening, native-CAD verification, and validation preparation. Unit-aware
engineering definitions and screening models were completed for all three
parts; the sampling and multi-start searches were executed; invalid
evaluations were contained; native geometry, physical properties, exchange
artifacts, and boundary evidence were produced; and analytical-to-CAD mass
agreement was assessed. Table~\ref{tab:objective-achievement} records the
status of each specific objective.

The aim was only partly achieved at the structural-validation level. The
bracket reached manual FEA review, but the two exploratory studies did not use
a sufficiently traceable and correct design load. Corrected, mesh-converged
FEA has not yet been completed for the selected bracket, padeye, or stabilizer
candidates.

\section{Limitations}
\label{sec:conclusion-limitations}

The analytical models are first-order representations of selected load paths.
They omit detailed fastener and weld behaviour, nonlinear contact,
three-dimensional local stress, fatigue, buckling, corrosion, vibration,
impact, and other component-specific service effects. The stabilizer load is a
prescribed structural envelope rather than a pressure field derived from CFD,
model tests, or measurements.

The search uses a finite 60-point sample, three local starts, fixed bounds, and
a penalty objective. A different seed, model, bound, or constraint strategy
could produce a different finalist. Several selected parameters lie near
their lower bounds, which indicates that manufacturing limits and omitted
failure modes can materially affect the result.

The CAD sweeps cover seeded design vectors rather than every point in the
continuous space. They verify geometry, mass, artifacts, and boundary
identification, but do not verify drawings, tolerances, assembly, weld access,
or manufacturability. No acceptable finalist FEA or physical test result is
available.

\section{Recommendations}
\label{sec:recommendations}

\begin{enumerate}
    \item Recreate the bracket study with a verified total load of 1472~N,
    confirm its coordinate direction, keep force per entity disabled, and
    demonstrate reaction balance. Compare the baseline, lightest finalist, and
    conservative backup candidates using at least three systematically refined
    meshes.
    \item Replace the rigid hole constraint with more realistic bolt, bearing,
    contact, and preload representations when these effects influence the
    bracket decision. Assess local plate behaviour, rib-root and weld stress,
    fatigue, corrosion allowance, fabrication tolerance, and supporting
    structure.
    \item Validate the padeye under the governing lifting standard and load
    combinations. Begin with a distributed bearing model and progress to a
    separate pin with clearance and nonlinear contact where required. Check
    net section, bearing, shear-out, lug root, gussets, welds, fatigue,
    inspection, and proof-load requirements.
    \item Derive stabilizer pressure cases from an accepted hydrodynamic method,
    CFD, model tests, or measurements. Include positive and negative loading,
    torsion, root and spar stresses, skin-panel buckling, fatigue, corrosion,
    actuator interaction, vibration, and tip displacement.
    \item Compare corrected FEA responses with the analytical predictions at
    identical candidate geometries and loads. Any calibration should use
    separate fitting and checking candidates and remain limited to the
    supported parameter range.
    \item Treat the lightest analytical row as one candidate rather than the
    automatic design choice. Material certification, applicable codes,
    uncertainty, fabrication, inspection, physical testing where required,
    and qualified engineering review must precede acceptance.
\end{enumerate}

\section{Future Work}
\label{sec:future-research}

Future development can introduce a formal multi-fidelity strategy in which
selected FEA results correct analytical bias or support a surrogate model with
explicit uncertainty. Reliability-based optimization can then represent
variability in loads, material properties, fabrication dimensions, corrosion,
and modelling error.

The search can also be extended from single-objective mass minimization to a
multi-objective formulation that includes fabrication cost, fatigue damage,
displacement, or embodied carbon. Manufacturing-aware constraints should
represent minimum gauges, available stock sizes, tool and inspection access,
weld size, bend radii, and tolerance sensitivity.

For the individual components, the bracket can be extended to equipment
vibration and foundation flexibility, the padeye to lifting-operation load
angles and proof testing, and the stabilizer to coupled hydrodynamic and
structural loading. Each extension should preserve the distinction between a
screening prediction, a CAD property, a solver result, and a validated design.
